% ===========================================================================
%  TJMJ — 桃江经典王麻将  全栈技术手册
%  Compile: xelatex TJMJ_Technical_Manual.tex (requires XeLaTeX + Chinese fonts)
% ===========================================================================
\documentclass[11pt,a4paper]{article}

% ---------- Font & encoding ----------
\usepackage[UTF8]{ctex}
\usepackage{fontspec}
\setmainfont{Times New Roman}
\setCJKmainfont{SimSun}[BoldFont=SimHei,ItalicFont=KaiTi]
\setCJKmonofont{FangSong}
\usepackage[scaled]{helvet}
\renewcommand{\ttdefault}{lmtt}

% ---------- Page layout ----------
\usepackage{geometry}
\geometry{top=2.5cm, bottom=2.5cm, left=2.5cm, right=2.5cm}
\usepackage{fancyhdr}
\pagestyle{fancy}
\fancyhf{}
\fancyhead[L]{\small\textit{TJMJ Technical Manual}}
\fancyhead[R]{\small\thepage}
\renewcommand{\headrulewidth}{0.4pt}

% ---------- Scientific packages ----------
\usepackage{amsmath,amssymb}
\usepackage{graphicx}
\usepackage[table]{xcolor}
\usepackage{booktabs}
\usepackage{tabularx}
\usepackage{enumitem}
\usepackage{hyperref}
\hypersetup{colorlinks=true,linkcolor=blue,citecolor=blue,urlcolor=blue}
\usepackage{float}
\usepackage{caption}
\usepackage{listings}
\usepackage{tikz}
\usetikzlibrary{shapes,arrows,positioning,fit}

% ---------- Code listings ----------
\lstset{
  basicstyle=\ttfamily\small,
  breaklines=true,
  frame=single,
  rulecolor=\color{gray!40},
  backgroundcolor=\color{gray!5},
  keywordstyle=\color{blue}\bfseries,
  commentstyle=\color{green!50!black},
  stringstyle=\color{red!70!black},
  showstringspaces=false,
  numbers=left,
  numberstyle=\tiny\color{gray},
  numbersep=5pt,
  tabsize=2,
  language={}
}

% ---------- Title / authors ----------
\title{\bfseries\LARGE TJMJ: A Full-Stack Web-Based Taojiang Classic \\[2pt]
       Wang Mahjong Platform for Entertainment \& Education \\[4pt]
       \large 全栈技术手册 v2.0}

\author{%
  JaiJaiC\thanks{Correspondence: \url{https://github.com/JaijaiC/TJMJ}}
  \and
  TJMJ Development Team \\
  \small\textit{Taojiang, Hunan, China} \\[4pt]
  \small\texttt{\{jaijai.c\}@tjmj.dev}
}
\date{\today}

% ---------- Abstract ----------
\newcommand{\natureabstract}[1]{%
  \begin{center}
  \rule{\linewidth}{0.8pt}\\[2pt]
  \begin{minipage}{0.95\linewidth}
  \small\textbf{Abstract}\quad #1
  \end{minipage}\\[2pt]
  \rule{\linewidth}{0.4pt}
  \end{center}
  \vspace{4pt}
}

\begin{document}
\maketitle

\natureabstract{
  桃江经典王麻将 (Taojiang Jingdian Wang Majiang, TJMJ) 是一款面向湖南桃江地区麻将爱好者的全栈网页麻将游戏平台。
  本项目完整实现了桃江地方麻将的核心规则——包括打筛扳王（癞子）机制、扎鸟翻倍、门子累加番、
  组合胡法等复杂计分体系，并提供单机模式、四人联机对战、观战模式三大玩法。
  技术架构采用 Vue 3 Composition API + Vite 8 前端、Node.js + ws WebSocket 服务端、
  WebRTC P2P 实时语音（含 Twilio TURN 中继），前端通过 GitHub Pages 托管、后端通过 ngrok 公网穿透。
  全文共八节，系统阐述游戏规则数学模型、前后端全栈架构（含完整通信协议与数据流）、
  核心算法实现（递归回溯胡牌判定、听牌检测、NPC 启发式策略）、
  多平台设备兼容性方案（iOS/Android/桌面端）、语音系统设计以及未来 Mahjong Agent 训练方向。
  本项目完全开源（MIT License），仅供家庭娱乐与学术研究使用。
}

\tableofcontents
\newpage

% ===================================================================
\section{Introduction}
% ===================================================================

麻将 (Mahjong) 作为中国传统博弈游戏的代表，在湖南桃江地区形成了独具特色的地方变种——桃江经典王麻将。
其核心特征包括：108 张万/筒/条字牌（不含风牌与花牌）、``王''（癞子/万能牌）机制、
打筛扳王确定癞子、扎鸟翻倍计分、门子累加番型等复杂规则，与标准国标麻将或四川血战麻将存在显著差异。

目前市面上仅有闲逸等少数 App 提供益阳麻将玩法，但其存在以下不足：
(1) 需要注册登录，操作繁琐；
(2) 规则与桃江本地玩法不完全一致；
(3) 缺乏拖拽手牌、实时语音、文字弹幕等功能；
(4) 无法训练 AI 对弈智能体。

为此，我们设计并开发了 TJMJ——一款与桃江麻将规则高度一致、无需安装、
点击链接即可进入的网页麻将平台。本文第 2 节详述游戏规则与数学模型；
第 3 节展示全栈系统架构与通信协议；第 4 节剖析核心算法实现；
第 5 节介绍联动功能（语音、头像、音效）；第 6 节分析多平台兼容性方案；
第 7 节总结已解决与仍存在的问题；第 8 节记录开发方法与部署流程。

\textbf{项目规模}：总计约 4200 行核心代码，其中前端主文件 \texttt{App.vue} 2600 行、
服务端 \texttt{GameRoom.js} 494 行、WebSocket 客户端 224 行、AI 策略 227 行、
胡牌引擎 208 行、规则检查 142 行、语音模块 100 行。

% ===================================================================
\section{Game Rules \& Mathematical Formulation}
% ===================================================================

\subsection{Tile Set \& Terminology}

桃江麻将使用精简牌组——共 108 张牌，三种花色各 36 张：

\begin{itemize}[nosep]
  \item \textbf{万} (Wan, $W$): 一万 $\sim$ 九万，各 4 张 (11--19)
  \item \textbf{条} (Tiao, $T$): 一条 $\sim$ 九条，各 4 张 (21--29)
  \item \textbf{筒} (Tong, $B$): 一筒 $\sim$ 九筒，各 4 张 (31--39)
\end{itemize}

编码约定：花色十位数（1=万, 2=条, 3=筒）+ 数字个位数。例：$38$ 表示八筒，$15$ 表示五万。

\textbf{术语表}：
\begin{itemize}[nosep]
  \item \textbf{王 (\textit{Wang})}： 癞子/万能牌，可替代任意牌使用。由开局``打筛扳王''确定。
  \item \textbf{地 (\textit{Di})}： 扳开的那张牌本身，用于判定地胡。
  \item \textbf{将 (\textit{Jiang})}： 数字为 2、5、8 的牌（任意花色），共 36 张。胡牌必需品。
  \item \textbf{一句话 (\textit{YiJuHua})}： 三张牌构成 $AAA$（刻子）或 $ABC$（顺子，同花色连续）。
  \item \textbf{听牌 (\textit{Ting})}： 13 张手牌差任意一张即可胡牌的状态。
  \item \textbf{扎鸟 (\textit{ZhaNiao})}： 胡牌后翻下一张牌，根据数字确定输赢翻倍。
  \item \textbf{门子 (\textit{MenZi})}： 特殊胡牌类型（碰碰胡、清一色等），采用累加番规则。
\end{itemize}

\subsection{Dice Rolling \& Wall Breaking}

开局由庄家掷两枚骰子，骰子点数记为 $d_1, d_2$（$d_1 \leq d_2$，$d_i \in \{1,\dots,6\}$）。
设 $S = d_1 + d_2$。

\textbf{定庄与数墙}：
庄家位置编号为 $P_0$，逆时针依次为下家 $P_1$、对家 $P_2$、上家 $P_3$。
从庄家开始数到第 $S$ 家（庄家$=1,5,9$），该玩家面前的牌墙即为起手墙。

\textbf{开牌位置}：在该墙右手边开始往左数，跳过 $d_1$（较小骰数）叠，从下一叠开始抓牌。

\textbf{抓牌顺序}：庄家起抓，每人每次取一码（两叠=4张），共三轮（$3 \times 4 = 12$ 张/人）。
然后``收庄''每人各补 1 张（各得 13 张），庄家再补 1 张（14 张），游戏开始。

\textbf{扳王规则}：
从起手墙对应玩家的下家墙左手边开始往右数，到第 $d_2$（较大骰数）叠处扳开最上面一张牌，
此牌即为``地''（$T_d$），$T_d + 1$（模 10 循环，花色不变）即为``王''/癞子：
\begin{equation}
  T_w = \begin{cases}
    \lfloor T_d/10 \rfloor \times 10 + 1, & \text{if } T_d \bmod 10 = 9 \\
    T_d + 1, & \text{otherwise}
  \end{cases}
\end{equation}

\subsection{Basic Actions \& Priority}

玩家回合可执行以下操作，优先级由高到低：

\begin{enumerate}[nosep,label=\textbf{\arabic*.}]
  \item \textbf{胡 (Hu)} — 点炮/自摸，直接获胜（最高优先级）
  \item \textbf{碰 (Pong)} / \textbf{杠 (Gang)} — 三张相同牌碰之；四张相同可开杠
  \item \textbf{吃 (Chi)} — 仅能吃上家打出的牌，构成顺子
\end{enumerate}

优先级规则：\textbf{点炮（胡）$>$ 碰（杠）$>$ 吃}。联机模式中，若同时可碰和可吃，两个按钮同时亮起供玩家选择。
抢杠规则：若明杠的牌恰好是别家要胡的牌，直接判杠家全赔。

\subsubsection*{吃 (Chi)}
只能吃上家打出的牌，且必须与手中两张牌构成同花色顺子。例：上家打 $4$ 万，手中有 $(2,3)$ 或 $(3,5)$ 或 $(5,6)$ 万即可吃。

\subsubsection*{碰 (Pong)}
任意玩家打出一张牌，若手中有两张相同即可碰。碰后三张亮出（副露），手牌数减 2。

\subsubsection*{杠 (Gang)}
分明杠（他人打出，手中有三张）和暗杠（手中四张，未亮出）。开杠后重新掷骰子，从牌墙末尾倒数摸三张。若摸到能胡的牌直接胡（杠上花），否则由下家开始依次判定能否胡此三张。

\subsection{Winning Patterns (Hu Types)}

胡牌的基本结构为：\textbf{一对将 (Jiang) + 四句话 (Sentence)}，
其中``一句话''可以是刻子（$AAA$）或顺子（$ABC$）。

表~\ref{tab:hu_types} 列出了所有胡法类型及其番数。门子胡法采用\textbf{累加番}规则。

\begin{table}[H]
\centering
\caption{胡法类型与番数对照表}
\label{tab:hu_types}
\small
\begin{tabular}{@{}llr@{}}
\toprule
\textbf{类型} & \textbf{条件} & \textbf{番数} \\
\midrule
平胡       & 有王，基本构成                & 1 番 \\
硬庄       & 无王，基本构成                & 2 番 \\
碰碰胡     & 四句话全为刻子 $AAA$          & 2 番 \\
清一色     & 手牌全部同花色               & 2 番 \\
七小对     & 14 张组成 7 个对子            & 2 番 \\
将将胡     & 14 张全部为 2/5/8             & 2 番 \\
起手胡     & 庄家 14 张首轮直接胡          & 2 番 \\
开杠       & 明杠/暗杠后摸三张胡牌         & 2 番 \\
地胡       & 摸到三张``地''即可胡          & 2 番 \\
天胡       & 手牌有 3 个王               & 3 番 \\
天天胡     & 手牌有 4 个王               & 5 番 \\
黑天胡     & 首轮无王、无将、无一句话      & 2 番 \\
地胡带拖   & 三张地不胡，当一句话用        & 3 番 \\
\bottomrule
\end{tabular}
\end{table}

\subsubsection*{组合胡法}
当多种门子同时满足时，番数累乘。通用计算公式：
\begin{equation}
  S_{\text{自摸}} = \left(\sum \text{门子倍数}\right) \times F_{\text{扎鸟}} \times F_{\text{杠}}
\end{equation}

实例：
\begin{itemize}[nosep]
  \item 天胡七小对全中：$(3+2) \times 2 = 10$ 番
  \item 天胡杠上花全中：$3 \times 2 \times 2 = 12$ 番
  \item 硬庄将将胡碰碰胡天胡杠上花全中：$(2+2+2+2) \times 2 \times 2 = 32$ 番
\end{itemize}

\subsection{Scoring \& ZhaNiao (扎鸟)}

胡牌分为\textbf{自摸}（自己摸到胡牌）和\textbf{抓炮}（他人打出胡牌）两种方式。

\textbf{自摸计分：} 胡牌后进行扎鸟：翻开本该下家摸的那张牌（$T_z$），根据 $T_z \bmod 10$ 判定：
\begin{itemize}[nosep]
  \item $1,5,9$（全中）：三家输分均翻倍
  \item $2,6$：下家输分翻倍
  \item $3,7$：对家输分翻倍
  \item $4,8$：上家输分翻倍
\end{itemize}

设基础分 $B = 3$（平胡 1 番 $\times$ 底分 3），番数 $M$，扎鸟系数 $Z \in \{1,2\}$：
赢家得分 $=$ 输家各支付 $B \times M \times Z$（扎中者）或 $B \times M$（未扎中者）。

\textbf{核心规则：} 有王只能自摸（不可抓炮），无王可以抓炮。这一规则使``硬庄''（无王胡牌）更具策略价值。

% ===================================================================
\section{System Architecture}
% ===================================================================

\subsection{Overview}

TJMJ 采用三层 B/S 架构。前端 Vue 3 SPA 部署于 GitHub Pages，服务端 Node.js WebSocket 运行于本地并通过 ngrok 暴露公网端点，WebRTC 语音为纯 P2P 网状拓扑。

\begin{figure}[H]
\centering
\begin{tikzpicture}[
  node distance=1.3cm,
  box/.style={rectangle, draw, rounded corners, minimum width=2.5cm, minimum height=0.9cm, align=center, font=\small},
  arrow/.style={->, >=stealth, thick}
]
  \node[box, fill=blue!10] (browser) {Browser Client\\Vue 3 + Vite\\TCP 5173};
  \node[box, fill=green!10, right=of browser] (ngrok) {ngrok Tunnel\\wss://};
  \node[box, fill=orange!10, right=of ngrok] (ws) {WebSocket Server\\Node.js :8080};
  \node[box, fill=red!10, below=of ws] (game) {Game Room FSM\\GameRoom.js};
  \node[box, fill=purple!10, below=of browser] (ai) {AI Agent\\NpcStrategy.js};
  \node[box, fill=cyan!10, left=of ai] (voice) {WebRTC Voice\\(P2P Mesh)};

  \draw[arrow] (browser) -- (ngrok) node[midway, above, font=\footnotesize] {wss://};
  \draw[arrow] (ngrok) -- (ws) node[midway, above, font=\footnotesize] {:8080};
  \draw[arrow] (ws) -- (game);
  \draw[arrow] (game) -- (ai);
  \draw[arrow] (browser) to[out=-60, in=60] node[right, font=\footnotesize] {P2P} (voice);
  \draw[arrow] (browser.east) to[out=30, in=150] (browser.west);
  \draw[arrow, dashed, gray] (ws) -- (voice) node[midway, right, font=\footnotesize] {信令中继};
\end{tikzpicture}
\caption{系统架构拓扑图}
\label{fig:arch}
\end{figure}

\subsection{Technology Stack}

\begin{table}[H]
\centering
\caption{技术栈总览}
\label{tab:stack}
\small
\begin{tabular}{@{}lll@{}}
\toprule
\textbf{层} & \textbf{技术} & \textbf{作用} \\
\midrule
前端框架 & Vue 3.5 (Composition API) & 响应式 UI \\
构建工具 & Vite 8 + Rolldown & 打包/HMR \\
前端部署 & GitHub Pages & 静态托管 \\
服务端运行时 & Node.js & 无框架纯 JS 服务 \\
服务端通信 & ws (v8.18) & WebSocket \\
服务端部署 & ngrok & TCP 隧道 (公网) \\
实时语音 & WebRTC (P2P 网状) & 4 人通话 \\
NAT 穿透 & Twilio TURN (10GB/月) & 移动网络穿透 \\
语音合成 & Web Speech API + 预录 mp3 & 牌名播报 \\
CI/CD & GitHub Actions & PR build 检查 \\
\bottomrule
\end{tabular}
\end{table}

\subsection{Project Structure}

完整的项目文件树及其职责：

\begin{lstlisting}[language=,basicstyle=\ttfamily\scriptsize]
TJMJ/
  frontend/                        # Vue 3 SPA
    src/
      App.vue                      # ★ 主组件 (2602行: UI+逻辑+样式)
      main.js                      # Vue 入口: createApp(App).mount('#app')
      style.css                    # 全局样式（iOS 安全区等）
      core/                        # 游戏引擎（纯逻辑，无 UI 依赖）
        constants.js               # 花色枚举、258 判定、王计算
        HuCalculator.js            # 递归回溯胡牌判定引擎 (208行)
        RuleChecker.js             # 吃/碰/杠/听牌合法性 (142行)
      ai/
        NpcStrategy.js             # NPC 启发式出牌策略 (227行)
      network/
        client.js                  # WebSocket 客户端 (224行: 重连/心跳/节流)
      utils/
        mjLogic.js                 # 洗牌算法 + 牌墙布局
        speech.js                  # 语音播报 (预录音优先 → TTS 降级)
      components/
        HelloWorld.vue             # Vite 脚手架模板（未使用）
    public/
      images/                      # 牌面素材(1-9万条筒 + 3D背景)
      sfx/                         # 预录制音效 (.gitkeep 预留)
      docs/                        # PDF教程 + 演示视频
      *.mp3                        # BGM 背景音乐 (1-13.mp3)
    index.html
    vite.config.js                 # base: '/TJMJ/' (GitHub Pages)
    package.json                   # vue 3.5 + vite 8.0 + gh-pages 6
  server/                          # Node.js 游戏服务器
    index.js                       # ★ WebSocket 入口 (247行)
    GameRoom.js                    # ★ 房间状态机 (494行)
    game/                          # 游戏逻辑（前端镜像，服务端权威校验）
      constants.js / HuCalculator.js / RuleChecker.js / mjLogic.js
    package.json                   # ws + uuid
  backend/                         # 早期 Socket.io 版本（已废弃）
  other/                           # 本地工具 (ngrok.exe + 演示视频)
  .github/workflows/deploy.yml     # GitHub Actions CI (build 检查)
  start-all.bat                    # 一键启动 (3 窗口)
  readme.md / ARCHITECTURE.md / .gitignore
\end{lstlisting}

\textbf{代码量统计：}

\begin{table}[H]
\centering
\caption{核心文件代码行数}
\label{tab:loc}
\small
\begin{tabular}{@{}lr@{}}
\toprule
\textbf{文件} & \textbf{行数} \\
\midrule
frontend/src/App.vue & 2602 \\
server/GameRoom.js & 494 \\
server/index.js & 247 \\
frontend/src/ai/NpcStrategy.js & 227 \\
frontend/src/network/client.js & 224 \\
frontend/src/core/HuCalculator.js & 208 \\
frontend/src/core/RuleChecker.js & 142 \\
frontend/src/utils/speech.js & 100 \\
\midrule
\textbf{合计} & \textbf{4244} \\
\bottomrule
\end{tabular}
\end{table}

\subsection{Communication Protocol}

客户端与服务器通过单一 WebSocket 长连接进行 JSON 消息通信。完整协议定义如下：

\begin{table}[H]
\centering
\caption{WebSocket 消息协议（完整版）}
\label{tab:protocol}
\footnotesize
\begin{tabular}{@{}llp{5.5cm}@{}}
\toprule
\textbf{type} & \textbf{方向} & \textbf{说明} \\
\midrule
\texttt{create\_room} & C$\to$S & 创建新房间（可选 4 位指定号码） \\
\texttt{join\_room} & C$\to$S & 加入房间（8888=测试） \\
\texttt{ready} & C$\to$S & 准备就绪（4 人满自动开局） \\
\texttt{ready\_next} & C$\to$S & 确认继续（4 人全确认开新局） \\
\texttt{discard} & C$\to$S & 出牌（携带牌值 tile） \\
\texttt{draw} & C$\to$S & 摸牌 \\
\texttt{action} & C$\to$S & 吃/碰/杠/胡操作（携带 action + combo） \\
\texttt{pass} & C$\to$S & 过牌 \\
\texttt{chat} & C$\leftrightarrow$S & 文字弹幕（携带 text + fromName） \\
\texttt{emoji} & C$\leftrightarrow$S & 表情（携带 icon + target） \\
\texttt{webrtc\_offer/answer/ice} & C$\leftrightarrow$S & WebRTC 信令中继（不解析 data） \\
\texttt{update\_avatar} & C$\to$S & 头像更新（携带 avatar 数据） \\
\texttt{get\_turn} & C$\to$S & 请求 TURN 凭据 \\
\texttt{join\_spectate} & C$\to$S & 加入观战 \\
\texttt{ping} & C$\to$S & 心跳检测 \\
\midrule
\texttt{room\_created/joined} & S$\to$C & 房间创建/加入成功 \\
\texttt{room\_players} & S$\to$C & 玩家列表更新（含昵称/头像/状态） \\
\texttt{game\_start} & S$\to$C & 开局（每人收到自己的 hand、wangTile、dice） \\
\texttt{game\_state} & S$\to$C & 实时状态同步（弃牌/副露/剩余牌数/当前回合） \\
\texttt{action\_prompt} & S$\to$C & 请求玩家操作（含 canHu/Peng/Gang/Chi） \\
\texttt{drew\_tile} & S$\to$C & 摸到某张牌 \\
\texttt{round\_end} & S$\to$C & 本局结束（含四家亮牌） \\
\texttt{pong} & S$\to$C & 心跳回复 \\
\texttt{turn\_credentials} & S$\to$C & Twilio TURN 凭据（动态生成） \\
\texttt{avatar\_updated} & S$\to$C & 头像已同步 \\
\texttt{game\_end} & S$\to$C & 16 局终局 \\
\bottomrule
\end{tabular}
\end{table}

\subsection{Server Architecture: GameRoom Finite State Machine}

\texttt{GameRoom.js} 是整个联机游戏的核心——它是一个有限状态机，管理从等待到结算的完整生命周期。

\textbf{状态转换：}
\begin{verbatim}
  LOBBY → (4人满 + 全部 ready) → PLAYING
  PLAYING → (有人胡牌 / 流局) → SETTLEMENT
  SETTLEMENT → (4人全部确认继续) → PLAYING (下一局)
  SETTLEMENT → (第16局结束) → LOBBY
\end{verbatim}

\textbf{核心方法：}

\begin{table}[H]
\centering
\caption{GameRoom 核心方法}
\label{tab:gameroom}
\small
\begin{tabular}{@{}ll@{}}
\toprule
\textbf{方法} & \textbf{职责} \\
\midrule
\texttt{startGame()} & 洗牌、发牌 (53张)、扳王、广播 game\_start \\
\texttt{handleDiscard(p, tile)} & 出牌处理：移除手牌 → 广播 → checkIntercepts \\
\texttt{handleDraw(p)} & 摸牌处理 → 自摸检测 \\
\texttt{checkIntercepts(p, tile)} & 检查 3 家能否吃碰杠胡 → 通知相关玩家 \\
\texttt{resolveActions()} & 收集响应 → 优先级排序（胡>杠>碰>吃）→ 执行 \\
\texttt{nextTurn()} & 切换回合 → 发新牌 → 通知摸牌 \\
\texttt{endRound(winner)} & 结算：赢家坐庄、亮牌、等待确认 \\
\texttt{playerReadyForNext(p)} & 确认计数，4 人满后开新局 \\
\texttt{broadcastPublic()} & 向所有玩家发送 game\_state \\
\bottomrule
\end{tabular}
\end{table}

\textbf{关键设计：}
\begin{itemize}[nosep]
  \item \textbf{双端校验}：frontend/core/ 和 server/game/ 各有一份相同的游戏逻辑（HuCalculator, RuleChecker, mjLogic）。单机模式走前端，联机模式前端仅做本地预览、服务端做权威判定。
  \item \textbf{超时保护}：\texttt{checkIntercepts} 设置 \texttt{pendingActions} 后会启动 10 秒超时。若某玩家断线未响应，超时自动强制清除并调用 \texttt{nextTurn()}，防止游戏永久卡死。
  \item \textbf{手牌不排序}：服务端发牌和摸牌均不调用 \texttt{.sort()}，保留发牌自然顺序。客户端通过计数差集算法智能合并——仅增删变更的牌，保持玩家手动拖拽顺序不变。
  \item \textbf{防重复抓牌}：\texttt{physicalDraw()} 自动跳过 \texttt{diIndex}（地牌位置），确保扳开的牌不会被摸走。
\end{itemize}

\subsection{Client Network Layer}

\texttt{client.js} 实现了生产级的 WebSocket 网络层：

\begin{itemize}[nosep]
  \item \textbf{自动重连}：指数退避策略 (1s → 2s → 4s → ... → 128s)，最多 8 次。重连成功后自动重新加入原房间（利用 URL 参数 \texttt{?auto\_join=} 保存房间号和昵称）。
  \item \textbf{心跳检测}：每 10 秒发 ping，收到 pong 后计算往返延迟，更新游戏界面左上角 4 格信号图标（<100ms 满格, <400ms 两格）。
  \item \textbf{消息节流}：游戏操作 400ms 节流防止高频崩溃；WebRTC 信令和聊天不节流。
  \item \textbf{缓冲队列}：连接建立前的消息先入队列，\texttt{onopen} 后批量冲刷。
  \item \textbf{URL 持久化}：创建/加入房间后自动将房间号和昵称写入 URL 参数，支持页面刷新后自动重连回当前对局。
\end{itemize}

\subsection{Deployment Topology}

\begin{verbatim}
  开发者本地                      云端
  ┌──────────┐   git push        ┌─────────────────┐
  │ 源码      │ ───────────────► │ GitHub (main)    │
  │ git repo  │                   │  ├── GitHub Pages│ ← 前端静态文件
  └────┬─────┘                   │  └── Actions CI  │ ← build 检查
       │ npm run deploy          └─────────────────┘
       │ (gh-pages -d dist)
       │
       ├── start-all.bat
       │   ├── Window 1: server → :8080
       │   ├── Window 2: frontend → :5173
       │   └── Window 3: ngrok http 8080
       ▼
    玩家访问 https://jaijaic.github.io/TJMJ/
         └── WebSocket → ngrok → server :8080
         └── WebRTC Voice → P2P mesh + Twilio TURN (10GB/月)
\end{verbatim}

% ===================================================================
\section{Core Algorithms}
% ===================================================================

\subsection{Hu Detection Algorithm (胡牌判定)}

胡牌判定是麻将引擎最核心的算法模块，对应文件 \texttt{HuCalculator.js}（前端和 server/game/ 各一份镜像）。
算法采用\textbf{递归回溯 + 贪心剪枝}策略：

\begin{enumerate}[nosep]
  \item \textbf{预处理}：分离王牌与普通牌，统计 $n_w$（王数量）和 $n_d$（地数量）
  \item \textbf{特殊拦截}：首轮 14 张检查黑天胡（无王、无将、无一句话）
  \item \textbf{顶级判定}：4 王→天天胡, 3 王→天胡, 3 地→地胡
  \item \textbf{常规判定}：调用 \texttt{checkNormalHu()} 递归去除刻子/顺子
  \item \textbf{七小对判定}：调用 \texttt{check7Pairs()} 贪心配对
  \item \textbf{门子判定}：将将胡→清一色→碰碰胡→硬庄→软庄，依序匹配
\end{enumerate}

\texttt{checkNormalHu(normalTiles, wangCount)} 的递归逻辑：
\begin{enumerate}[nosep]
  \item 若王充足，补齐对子（将）后递归去除句子
  \item 每次选择最小的牌 $t_{\min}$，尝试三种消除：
    \begin{itemize}[nosep]
      \item 刻子：$\{t,t,t\}$ 三张相同
      \item 顺子：$\{t,t+1,t+2\}$ 三张连续（同花色）
      \item 对子：$\{t,t\}$ 作为将（仅当将未确定时）
    \end{itemize}
  \item 若普通牌无法匹配，消耗 1 个王替代，继续递归
  \item 全部消除成功 → 胡牌；任何分支失败 → 回溯
\end{enumerate}

计算复杂度：最坏情况 $O(3^k)$，$k$ 为递归深度，通常 $k \leq 4$，单次判定 <1ms。

\subsection{Ting Detection Algorithm (听牌检测)}

听牌检测用于判断 13 张手牌是否差一张即可胡牌，以及开杠后是否能听牌。
对应 \texttt{RuleChecker.isTing()}：

\begin{enumerate}[nosep]
  \item 手牌数若不为 $3n+1$，直接返回 false
  \item 枚举所有可能的 34 种牌（效率优化：仅搜索手牌已有花色 $\pm 2$ 范围）
  \item 对每种候选牌 $t$，模拟加入手牌后调用 \texttt{HuCalculator.checkHu()}
  \item 若任意候选能胡 → 听牌 true；否则 false
\end{enumerate}

计算复杂度 $O(34 \times 3^k)$，实战 <1ms。

\subsection{NPC Strategy Agent (AI 出牌决策)}

\texttt{NpcStrategy.js} (227 行) 实现启发式决策：

\begin{enumerate}[nosep]
  \item \textbf{听牌优先}：遍历每张手牌，模拟丢弃后检查是否听牌 (\texttt{findBestDiscardToTing()})。若能直接听牌，优先打出该牌。
  \item \textbf{价值排序}：若未听牌，按优先级评估手牌价值 (\texttt{chooseWorstTile()})：
    \begin{itemize}[nosep]
      \item 孤张（无相邻牌、非将、非王）→ 最优先丢弃
      \item 边张/坎张（只能等 1 张进张）→ 次优先
      \item 搭子（差 1 张成顺子/刻子）→ 保留
      \item 将对（2/5/8 对子）→ 高价值保留
      \item 王牌 → 绝不丢弃
    \end{itemize}
  \item \textbf{安全出牌}：参考弃牌池，避免打出别家可能胡的炮牌
  \item \textbf{吃碰决策}：\texttt{shouldPeng()} / \texttt{shouldChi()} 评估手牌改善程度
\end{enumerate}

未来方向：升级为 DRL（深度强化学习）+ MCTS（蒙特卡洛树搜索）的自我对弈训练智能体。

\subsection{Voice System}

\texttt{speech.js} (100 行) 实现双轨语音播报：

\begin{enumerate}[nosep]
  \item \textbf{预录制优先}：检测 \texttt{public/sfx/} 下是否存在对应 mp3 文件（如 \texttt{tile\_15.mp3} → ``五万''），存在则直接播放 Audio 元素
  \item \textbf{TTS 降级}：文件缺失时使用浏览器 \texttt{SpeechSynthesis} API 合成语音
  \item \textbf{Web Audio 合成}：``咚''出牌音由 OscillatorNode 实时合成（800Hz → 200Hz 扫描，150ms）
\end{enumerate}

\textbf{音效文件清单}（需录制 34 个文件放入 \texttt{public/sfx/}）：
27 个牌名 (\texttt{tile\_11.mp3}--\texttt{tile\_39.mp3}) + 7 个动作词 (\texttt{chi.mp3}, \texttt{peng.mp3}, \texttt{gang.mp3}, \texttt{hu.mp3}, \texttt{zimo.mp3}, \texttt{win.mp3}, \texttt{pass.mp3})。

\subsection{WebRTC Voice Architecture}

联机语音采用\textbf{P2P 网状拓扑}（Mesh）：每个玩家与房间内其他 3 人各建立一条 RTCPeerConnection。
信令（offer/answer/ICE candidate）通过 WebSocket 服务端中继。

\begin{itemize}[nosep]
  \item \textbf{STUN/TURN}：Google STUN ×2 + Twilio TURN（10GB/月免费额度）。
  TURN 凭据由服务端 HMAC-SHA1 动态生成，24 小时有效，每小时刷新。
  \item \textbf{音频约束}：使用 \texttt{echoCancellation, noiseSuppression, autoGainControl} 三项降噪。
  \item \textbf{增量连接}：不因玩家列表变化而拆毁已有连接；仅在玩家加入/离开时增量创建/移除对应 PeerConnection。
  \item \textbf{音轨替换}：麦克风重新开启时，遍历已有连接替换 senders 音轨。
  \item \textbf{iOS 兼容}：AudioContext 创建后检测 \texttt{suspended} 状态并显式 \texttt{resume()}。
\end{itemize}

% ===================================================================
\section{Features \& User Interface}
% ===================================================================

\subsection{Game Modes}

平台提供三种游戏模式，均通过 \texttt{enterGame(mode)} 入口统一调度：

\begin{enumerate}[nosep]
  \item \textbf{单机模式}：玩家 vs 3 个 NPC AI。适合练习、测试规则。由客户端 \texttt{startRound()} 完全控制游戏流程。
  \item \textbf{联机模式}：4 人实时对战。创建/加入 4 位数字房间号。服务端 GameRoom 管理状态，客户端仅做本地预览。
  \item \textbf{观战模式}：旁观正在进行的对局。接收 \texttt{game\_state} 同步渲染，可选择观看任意一方视角。
\end{enumerate}

\textbf{模式切换数据流：}
\begin{verbatim}
  isInMenu=true
    ├── 点击"单机模式" → enterGame('single')
    │   └── handleReady() → startRound() → 本地游戏循环
    ├── 点击"联机模式" → enterGame('multi')
    │   └── createRoom() / joinRoom() → WebSocket → 等待4人 → 服务器开局
    ├── 点击"观战模式" → enterGame('spectate')
    │   └── joinSpectate() → 接收 game_state 并显示
    └── 点击"其他玩法" → 已禁用
\end{verbatim}

\subsection{Interactive Features}

\begin{itemize}[nosep]
  \item \textbf{拖拽手牌}：HTML5 Drag \& Drop + 触屏 Touch Events 双支持。手牌稳定 ID key 防止 Vue VNode 重排。所有游戏操作（出牌/吃/碰/杠）均同步更新 \texttt{handTileIds} 数组保持顺序一致。
  \item \textbf{头像系统}：点击自己头像弹出 12 格预设选择面板 + 拍照/上传自定义图片（FileReader → base64 data URL）。联机模式下通过 \texttt{update\_avatar} 消息实时同步给其他三人。头像存储于 \texttt{localStorage} 持久化。
  \item \textbf{表情互动}：5 种表情（🌹💩🥚💣💵）点击发送，飞行动画效果。💵 真转账 1 分。
  \item \textbf{文字弹幕}：右上角聊天按钮，消息从右向左滚动，兼容桌面和移动端。
  \item \textbf{BGM 系统}：首页 2 首交替播放 + 游戏内 9-13 首随机不重复队列 + O 叔钢琴模式（7 首）。统一音量压缩（Web Audio DynamicsCompressor）。
  \item \textbf{实时语音}：WebRTC P2P 网状通话，Twilio TURN 中继。音量实时显示条形波动。
  \item \textbf{语音播报}：出牌时自动播报牌名（``五万''``碰''等）。
\end{itemize}

% ===================================================================
\section{Device Compatibility}
% ===================================================================

\subsection{Mobile Device Handling}

本项目已实现以下跨平台兼容措施：

\begin{table}[H]
\centering
\caption{兼容性修复清单}
\label{tab:compat}
\small
\begin{tabular}{@{}ll@{}}
\toprule
\textbf{问题} & \textbf{解决方案} \\
\midrule
手机 NAT 穿透失败 & Twilio TURN (10GB/月) + Google STUN ×2 \\
iOS Safari click 不触发 & \texttt{<div>} → \texttt{<button>} 替换非交互元素 \\
iOS AudioContext 静音 & 检测 \texttt{suspended} 状态并 \texttt{resume()} \\
\multicolumn{1}{l}{iOS 双击缩放} & \texttt{@touchend} handler 中 \texttt{preventDefault()} \\
Android 地址栏跳动 & \texttt{100dvh} 替代 \texttt{100vh}（带 fallback） \\
触屏慢点击漏检 & 轻点阈值从 300ms 扩大到 500ms \\
WebSocket 断线 & 指数退避重连 (1s→128s, 最多8次) \\
桌面端 fixed 弹窗偏移 & 弹窗组件移出 \texttt{game-wrapper} 避免 transform 影响 \\
联机操作卡死 & 服务端 10 秒超时强制清除 \texttt{pendingActions} \\
麦克风权限拒绝 & 特征检测 + 用户提示 \\
\bottomrule
\end{tabular}
\end{table}

\subsection{Known Remaining Issues}

\begin{enumerate}[nosep]
  \item \textbf{语音偶发断续}：P2P 网状拓扑在 3+ 人时复杂度 $O(n^2)$，高丢包环境下部分连接可能掉线。长期方案：引入 SFU（Selective Forwarding Unit）中转。
  \item \textbf{桌面端 CSS 缩放}：\texttt{updateGameScale()} 对宽屏/竖屏适配仍有边界情况。
  \item \textbf{手牌拖拽范围}：触屏拖拽时 DOM 查询 \texttt{getBoundingClientRect()} 在 CSS Transform 缩放后可能返回不准确坐标。
  \item \textbf{音效文件缺失}：34 个预录制 mp3 文件尚未录制，当前全部降级为 TTS。
  \item \textbf{观战模式功能有限}：仅支持选择一方视角，不支持自动切换。
\end{enumerate}

% ===================================================================
\section{Development \& Deployment}
% ===================================================================

\subsection{Development Environment}

\begin{table}[H]
\centering
\caption{开发环境与依赖}
\label{tab:env}
\small
\begin{tabular}{@{}ll@{}}
\toprule
\textbf{组件} & \textbf{版本/技术} \\
\midrule
前端框架 & Vue 3.5 (Composition API + \texttt{<script setup>}) \\
构建工具 & Vite 8.0 + Rolldown \\
服务端 & Node.js + ws 8.18 \\
语音 & Web Speech API + WebRTC + 预录 mp3 \\
TURN 中继 & Twilio TURN (10GB/月，HMAC-SHA1 动态凭据) \\
内网穿透 & ngrok (前端 localhost 5173，后端 8080) \\
版本管理 & Git + GitHub \\
前端部署 & GitHub Pages (\texttt{gh-pages -d dist}) \\
CI & GitHub Actions (PR 触发 build 检查) \\
\midrule
Node 依赖 & vue, vite, @vitejs/plugin-vue, gh-pages \\
 & ws, uuid \\
\bottomrule
\end{tabular}
\end{table}

\subsection{Startup Commands}

\textbf{一键启动（Windows）：}
\begin{lstlisting}[language=bash, basicstyle=\ttfamily\small]
.\start-all.bat
# 依次启动 3 个窗口:
#   [1] Game Server   WebSocket :8080
#   [2] Frontend      Vite :5173
#   [3] ngrok Tunnel  (forward 8080, optional)
\end{lstlisting}

\textbf{手动启动：}
\begin{lstlisting}[language=bash, basicstyle=\ttfamily\small]
# Terminal 1 - Game Server
cd server && npm run dev

# Terminal 2 - Frontend
cd frontend && npm run dev

# Terminal 3 - ngrok (optional, for public access)
ngrok http 8080
\end{lstlisting}

\textbf{部署前端：}
\begin{lstlisting}[language=bash, basicstyle=\ttfamily\small]
cd frontend
npm run build      # Vite 打包到 dist/
npm run deploy     # gh-pages 推送至 GitHub Pages
\end{lstlisting}

\textbf{环境变量（TURN 功能需要）：}
\begin{lstlisting}[language=bash, basicstyle=\ttfamily\small]
set TWILIO_ACCOUNT_SID=ACxxxxxxxxxx
set TWILIO_AUTH_TOKEN=your_auth_token
cd server && npm run dev
\end{lstlisting}

\subsection{Data Flow: A Complete Game Round}

\begin{verbatim}
玩家 A 打开页面
  └── onMounted → restoreAvatar → BGM 默认关闭

玩家 A 创建房间
  └── connect() → create_room → server: new GameRoom
      └── A 收到 room_created (roomId)
      └── URL 写入 ?auto_join=xxxx&name=AAA

玩家 B/C/D 加入房间
  └── join_room(roomId) → server: addPlayer
      └── room_players 广播 → 4人列表显示

4人 ready → 服务器开局
  └── GameRoom.startGame()
      ├── Fisher-Yates 洗牌 108 张
      ├── 掷骰子 (d1,d2) → 定庄家起手位置
      ├── 扳王 → calculateWang(diTile)
      ├── 发牌 53 张 (庄家14张, 闲家13张)
      └── game_start 广播 (每人只收到自己的 hand)

游戏循环:
  A 出牌 → send({type:'discard',tile})
    → server: handleDiscard → 移除手牌 → broadcastPublic
    → 其他人收到 game_state: 更新弃牌池 + 手牌数

  checkIntercepts(A→B/C/D):
    ├── 有人能胡 → endRound(winner)
    ├── 有人能碰/杠 → action_prompt (等待选择)
    ├── 下家能吃 → action_prompt (等待选择)
    └── 无人拦截 → nextTurn → B 摸牌 → broadcastPublic

  B 摸牌 → 自摸检测 → 出牌 → 循环继续
  超时保护: 10s 无响应 → 强制跳过

胡牌 → endRound(winnerIndex)
  ├── 赢家坐庄 (dealerIndex = winnerIndex)
  ├── 亮牌 showdownHands (四家手牌全部可见)
  ├── 各玩家点击"确认继续" → playerReadyForNext
  └── 4人全部确认 → startGame (下一局)
\end{verbatim}

% ===================================================================
\section*{Future Work}
\addcontentsline{toc}{section}{Future Work}

\begin{itemize}[nosep]
  \item \textbf{Mahjong Agent}：基于深度强化学习（PPO/DQN）+ MCTS 自我对弈，训练达到人类高手水平的 AI。
  \item \textbf{语音升级}：引入 SFU 中转替代 P2P 网状，解决 4 人语音不稳定问题。
  \item \textbf{音效录制}：完成 34 个预录制 mp3 文件，提升游戏沉浸感。
  \item \textbf{观战增强}：支持多视角自动切换、回放功能。
  \item \textbf{性能优化}：App.vue 分拆为多个组件，降低单文件复杂度。
  \item \textbf{移动端原生体验}：PWA 离线支持、推送通知。
\end{itemize}

% ===================================================================
\section*{Code Availability}
\addcontentsline{toc}{section}{Code Availability}

本项目完全开源，代码仓库：\url{https://github.com/JaijaiC/TJMJ}

欢迎 Fork / Star / Pull Request。

% ===================================================================
\section*{Disclaimer}
\addcontentsline{toc}{section}{Disclaimer}

本项目为个人娱乐学习用途，不涉及任何商业行为。
所有规则及实现仅供参考，实际游玩请以当地线下规则为准。
因使用本软件产生的任何争议或损失，开发者不承担相关责任。

% ===================================================================
\section*{Acknowledgments}
\addcontentsline{toc}{section}{Acknowledgments}

感谢湖南桃江的家人朋友提供规则校验与测试反馈；感谢开源社区
(Vue.js, Node.js, Vite, ws) 提供的优秀工具链。

参考项目：
\begin{itemize}[nosep]
  \item OpenMajiang (\url{https://gitee.com/mirrors/OpenMajiang})
  \item 武汉麻将在线 (\url{http://cdn0001.afrxvk.cn/whmj/go.html})
\end{itemize}

% ===================================================================
\vfill
\begin{center}
\rule{\linewidth}{0.4pt}
{\small\textit{TJMJ v2.0 -- Technical Manual -- \today}}\\[2pt]
{\tiny 桃江经典王麻将\,\,|\, Taojiang Classic Wang Mahjong}\\[2pt]
{\tiny \copyright\ 2024--2026 JaiJaiC. All rights reserved. 仅供个人娱乐学习使用。}
\end{center}

\end{document}
